\begin{table}[H]
  \centering
  \small
  \caption{Synthèse des méthodologies existantes}
  \label{tab:synthese-methodes}
  \PFETableSetup
  \PFETableRowColors
  \PFETableArrayStretch
  \PFETableTabularXBegin
  \begin{tabularx}{\PFETableBlockWidth}{@{}%
    >{\raggedright\arraybackslash\hsize=0.65\hsize}X%
    >{\raggedright\arraybackslash\hsize=1.1166667\hsize}X%
    >{\raggedright\arraybackslash\hsize=1.1166667\hsize}X%
    >{\raggedright\arraybackslash\hsize=1.1166667\hsize}X%
  @{}}
    \tableHeaderStyle \tableHeaderCell{Méthodologie} & \tableHeaderCell{Scrum} & \tableHeaderCell{RUP} & \tableHeaderCell{GIMSI} \\
    Nombre d’étapes & 5 & 4 & 4 \\
    Étapes principales &
      Initialisation, Planification, Conception, Implémentation\newline
      Revue \& Rétrospective &
      Inception\newline
      Élaboration\newline
      Construction\newline
      Transition &
      Identification\newline
      Préparation\newline
      Implémentation\newline
      Évaluation \\
    Orientation & Développement logiciel agile & Systèmes d’information & BI \& SI \\
    Structure & Flexible et itérative & Itérative et structurée & Hautement structurée \\
    Flexibilité & Élevée & Moyenne & Faible \\
    Gestion des changements & Très bonne & Moyenne & Faible \\
    Adaptation au projet CSR & Très adaptée & Adaptée & Adaptée \\
  \end{tabularx}%
  \PFETableTabularXEnd
\end{table}
